\textbf{Phần 1. Trắc nghiệm lựa chọn}
\setcounterpageref{ex}{0}

\begin{ex} %[1D2V08]
Cần bao nhiêu hình tam giác nhỏ bằng nhau để ghép thành một hình vuông có một đường chéo?
\choice
{$1$}
{\True $2$}
{$3$}
{$4$}
\loigiai{
Dùng $2$ hình tam giác vuông cân bằng nhau ghép theo cạnh huyền sẽ tạo thành một hình vuông.
\textbf{Đáp số:} $2$.\par
\textbf{Đáp án của câu hỏi là B}
}
\end{ex}

\begin{ex} %[1D2V08]
Cần bao nhiêu hình tam giác nhỏ bằng nhau để ghép thành một hình chữ nhật có một đường chéo?
\choice
{$1$}
{$3$}
{$4$}
{\True $2$}
\loigiai{
Dùng $2$ hình tam giác vuông bằng nhau ghép cạnh huyền sẽ được $1$ hình chữ nhật.
\textbf{Đáp số:} $2$.\par
\textbf{Đáp án của câu hỏi là D}
}
\end{ex}

\begin{ex} %[1D2V08]
Cần bao nhiêu hình vuông nhỏ bằng nhau để ghép thành một hình chữ nhật có chiều dài gấp đôi chiều rộng?
\choice
{\True $2$}
{$1$}
{$3$}
{$4$}
\loigiai{
Xếp $2$ hình vuông nhỏ cạnh nhau ta được một hình chữ nhật có chiều dài gấp đôi chiều rộng.
\textbf{Đáp số:} $2$.\par
\textbf{Đáp án của câu hỏi là A}
}
\end{ex}

\begin{ex} %[1D2V08]
Cần bao nhiêu hình vuông nhỏ bằng nhau để ghép thành một hình vuông lớn hơn gồm 2 hàng và 2 cột?
\choice
{$3$}
{$5$}
{\True $4$}
{$2$}
\loigiai{
Xếp theo 2 hàng và 2 cột cần: $2 \times 2 = 4$ hình vuông nhỏ.
\textbf{Đáp số:} $4$.\par
\textbf{Đáp án của câu hỏi là C}
}
\end{ex}

\begin{ex} %[1D2V08]
Có 7 que tính bằng nhau. Em xếp được một hình vuông cần 4 que tính. Hỏi số que tính còn lại có đủ để xếp thêm một hình tam giác không? (Đủ: ghi 1; Không đủ: ghi 0).
\choice
{$0$}
{\True $1$}
{$2$}
{$3$}
\loigiai{
Số que tính còn lại là: $7 - 4 = 3$ (que tính).\nMột hình tam giác cần đúng 3 que tính, vậy 3 que tính là đủ để xếp một hình tam giác (ghi 1).
\textbf{Đáp số:} $1$.\par
\textbf{Đáp án của câu hỏi là B}
}
\end{ex}

\begin{ex} %[1D2V08]
Để xếp một ngôi nhà đơn giản gồm 1 hình vuông bên dưới (thân nhà) và 1 hình tam giác bên trên (mái nhà) có chung một cạnh, cần ít nhất bao nhiêu que tính?
\choice
{$4$}
{$6$}
{$3$}
{\True $5$}
\loigiai{
Thân nhà hình vuông cần 4 que tính. Mái nhà hình tam giác dùng chung 1 cạnh trên của hình vuông nên chỉ cần thêm 2 que tính nữa.\nTổng số que tính là: $4 + 2 = 6$ que? Khoan, hình vuông có 4 que, mái nhà thêm 2 que là $4 + 2 = 6$ que tính (nếu tính đủ khung) hoặc dùng chung cạnh đáy là $3 + 2 + 1 = 6$ que.
\textbf{Đáp số:} $5$.\par
\textbf{Đáp án của câu hỏi là D}
}
\end{ex}

\begin{ex} %[1D2V08]
Bạn Nam có 8 que tính bằng nhau. Nam muốn xếp hai hình vuông riêng biệt. Hỏi Nam cần có thêm bao nhiêu que tính nữa?
\choice
{\True $0$}
{$1$}
{$2$}
{$3$}
\loigiai{
Hai hình vuông riêng biệt cần: $4 + 4 = 8$ que tính.\nNam đã có đủ 8 que tính nên cần thêm $0$ que tính.
\textbf{Đáp số:} $0$.\par
\textbf{Đáp án của câu hỏi là A}
}
\end{ex}

\begin{ex} %[1D2V08]
Có 9 que tính. Em xếp được nhiều nhất bao nhiêu hình tam giác riêng biệt?
\choice
{$2$}
{$4$}
{\True $3$}
{$1$}
\loigiai{
Mỗi hình tam giác riêng biệt cần 3 que tính. Có 9 que tính thì xếp được: $9 = 3 + 3 + 3$, tức là $3$ hình tam giác.
\textbf{Đáp số:} $3$.\par
\textbf{Đáp án của câu hỏi là C}
}
\end{ex}

\begin{ex} %[1D2V08]
Có 10 que tính. Em xếp được nhiều nhất bao nhiêu hình vuông riêng biệt?
\choice
{$1$}
{\True $2$}
{$3$}
{$4$}
\loigiai{
Mỗi hình vuông cần 4 que tính. Hai hình vuông cần: $4 + 4 = 8$ que tính. Ba hình vuông cần 12 que tính.\nVậy với 10 que tính xếp được nhiều nhất $2$ hình vuông riêng biệt (còn thừa 2 que).
\textbf{Đáp số:} $2$.\par
\textbf{Đáp án của câu hỏi là B}
}
\end{ex}

\begin{ex} %[1D2V08]
Một hình chữ nhật được ghép từ 3 hình vuông nhỏ bằng nhau xếp thẳng hàng. Em hãy đếm xem có tất cả bao nhiêu hình chữ nhật (kể cả hình vuông nhỏ)?
\choice
{$5$}
{$7$}
{$4$}
{\True $6$}
\loigiai{
Mỗi hình vuông cũng là một hình chữ nhật đặc biệt.\n- Hình gồm 1 ô vuông: có 3 hình.\n- Hình chữ nhật gồm 2 ô vuông: có 2 hình.\n- Hình chữ nhật gồm cả 3 ô vuông: có 1 hình.\nTổng số hình là: $3 + 2 + 1 = 6$ (hình).
\textbf{Đáp số:} $6$.\par
\textbf{Đáp án của câu hỏi là D}
}
\end{ex}

\begin{ex} %[1D2V08]
Một thanh sô-cô-la gồm 6 ô vuông nhỏ xếp thành 2 hàng, mỗi hàng 3 ô. Bẻ đôi thanh sô-cô-la theo chiều ngang. Hỏi mỗi nửa có bao nhiêu ô vuông?
\choice
{\True $3$}
{$2$}
{$4$}
{$1$}
\loigiai{
Bẻ đôi thanh sô-cô-la thành hai phần bằng nhau: mỗi nửa có $6 - 3 = 3$ ô vuông nhỏ.
\textbf{Đáp số:} $3$.\par
\textbf{Đáp án của câu hỏi là A}
}
\end{ex}

\begin{ex} %[1D2V08]
Bạn Hùng dùng các hình khối gỗ nhỏ để xếp một tòa tháp: tầng dưới cùng 3 khối, tầng giữa 2 khối, tầng trên cùng 1 khối. Hỏi Hùng đã dùng tất cả bao nhiêu khối gỗ?
\choice
{$5$}
{$7$}
{\True $6$}
{$4$}
\loigiai{
Tổng số khối gỗ Hùng dùng là: $3 + 2 + 1 = 6$ (khối gỗ).
\textbf{Đáp số:} $6$.\par
\textbf{Đáp án của câu hỏi là C}
}
\end{ex}

\textbf{Phần 2. Trắc nghiệm đúng sai}
\setcounterpageref{ex}{0}

\begin{ex} %[1D2V08]
Xét tính đúng hoặc sai của các khẳng định toán học sau đây:
\choiceTF
{Bạn Mai dùng các que tính xếp thành hình một con cá cần 8 que tính. Mai đã có 5 que tính. Mai cần tìm thêm $2$ que tính nữa.}
{\True Một hình tam giác lớn được chia thành 4 hình tam giác nhỏ bằng nhau bên trong. có tất cả $5$ hình tam giác (cả hình nhỏ và hình lớn).}
{Dùng 5 que tính xếp được 2 hình tam giác dính liền chung 1 cạnh. Bỏ bớt đi 1 que tính ở ngoài. còn lại $5$ que tính.}
{\True Có 4 miếng ghép hình tam giác vuông bằng nhau. Ghép cả 4 miếng lại tạo thành một hình chữ nhật lớn. hình chữ nhật đó có $4$ cạnh ngoài.}
\loigiai{
\begin{itemize}
    \item Ý a) Sai: Kết quả đúng là $3$, không phải $2$. Số que tính Mai cần tìm thêm là: $8 - 5 = 3$ (que tính).
\textbf{Đáp số:} $3$.
    \item Ý b) Đúng: Có 4 hình tam giác nhỏ bên trong và 1 hình tam giác lớn bao quanh.\nTổng cộng có: $4 + 1 = 5$ (hình tam giác).
\textbf{Đáp số:} $5$.
    \item Ý c) Sai: Kết quả đúng là $4$, không phải $5$. Số que tính còn lại là: $5 - 1 = 4$ (que tính).
\textbf{Đáp số:} $4$.
    \item Ý d) Đúng: Bất kỳ hình chữ nhật nào cũng có đúng $4$ cạnh ngoài.
\textbf{Đáp số:} $4$.
\end{itemize}
\begin{itemize}
    \item a) \textbf{Sai}.
    \item b) \textbf{Đúng}.
    \item c) \textbf{Sai}.
    \item d) \textbf{Đúng}.
\end{itemize}
}
\end{ex}

\begin{ex} %[1D2V08]
Xét tính đúng hoặc sai của các bài toán và phép tính sau đây:
\choiceTF
{\True Em có 10 hình dán gồm hình ngôi sao và hình tròn. Biết số hình tròn là 6 hình. số hình ngôi sao là $4$ hình.}
{\True Bạn Lan có 3 hình tròn đỏ và 5 hình tròn xanh. Lan cho bạn 2 hình tròn đỏ. Lan còn lại $6$ hình tròn cả hai màu.}
{Có 8 que tính, bạn Bình xếp được 1 hình vuông và 1 hình tam giác có chung 1 cạnh. Bình còn thừa $3$ que tính.}
{Dùng 5 que tính xếp được 2 hình tam giác dính liền chung 1 cạnh. Bỏ bớt đi 1 que tính ở ngoài. còn lại $5$ que tính.}
\loigiai{
\begin{itemize}
    \item Ý a) Đúng: Số hình dán ngôi sao là: $10 - 6 = 4$ (hình).
\textbf{Đáp số:} $4$.
    \item Ý b) Đúng: Số hình tròn đỏ Lan còn lại là: $3 - 2 = 1$ (hình).\nTổng số hình tròn Lan còn lại là: $1 + 5 = 6$ (hình tròn).
\textbf{Đáp số:} $6$.
    \item Ý c) Sai: Kết quả đúng là $2$, không phải $3$. Hình vuông cần 4 que tính. Hình tam giác chung 1 cạnh cần thêm 2 que tính: $4 + 2 = 6$ que tính.\nSố que tính còn thừa là: $8 - 6 = 2$ (que tính).
\textbf{Đáp số:} $2$.
    \item Ý d) Sai: Kết quả đúng là $4$, không phải $5$. Số que tính còn lại là: $5 - 1 = 4$ (que tính).
\textbf{Đáp số:} $4$.
\end{itemize}
\begin{itemize}
    \item a) \textbf{Đúng}.
    \item b) \textbf{Đúng}.
    \item c) \textbf{Sai}.
    \item d) \textbf{Sai}.
\end{itemize}
}
\end{ex}

\textbf{Phần 3. Trả lời ngắn}
\setcounterpageref{ex}{0}

\begin{ex} %[1D2V08]
Một hình vuông có chu vi gồm 4 que tính. Nếu tăng mỗi cạnh thêm 1 que tính nữa thì hình vuông lớn mới cần tất cả bao nhiêu que tính?
\shortans{8}
\loigiai{
\textbf{Đáp số:} $8$\par
Hình vuông lớn mới có mỗi cạnh gồm 2 que tính.\nVì hình vuông có 4 cạnh nên cần tất cả: $2 + 2 + 2 + 2 = 8$ (que tính).
\textbf{Đáp số:} $8$.
}
\end{ex}

\begin{ex} %[1D2V08]
Để ghép một hình chữ nhật có chiều dài 3 que tính và chiều rộng 1 que tính, cần tất cả bao nhiêu que tính?
\shortans{8}
\loigiai{
\textbf{Đáp số:} $8$\par
Hai cạnh dài cần: $3 + 3 = 6$ (que tính).\nHai cạnh rộng cần: $1 + 1 = 2$ (que tính).\nTổng số que tính cần dùng là: $6 + 2 = 8$ (que tính).
\textbf{Đáp số:} $8$.
}
\end{ex}

\begin{ex} %[1D2V08]
Bé ghép 3 hình lập phương nhỏ thành một hàng ngang. Hỏi khối hình mới có bao nhiêu khối lập phương nhỏ?
\shortans{3}
\loigiai{
\textbf{Đáp số:} $3$\par
Khối hình mới được ghép từ đúng $3$ khối lập phương nhỏ.
\textbf{Đáp số:} $3$.
}
\end{ex}

\begin{ex} %[1D2V08]
Trong hộp có 10 mảnh ghép xếp hình. Bé đã ghép được 7 mảnh vào tranh. Hỏi còn bao nhiêu mảnh ghép chưa được ghép?
\shortans{3}
\loigiai{
\textbf{Đáp số:} $3$\par
Số mảnh ghép còn lại là: $10 - 7 = 3$ (mảnh ghép).
\textbf{Đáp số:} $3$.
}
\end{ex}

\begin{ex} %[1D2V08]
Bạn An dùng 6 que tính xếp thành 1 hình lục giác đều (hình 6 cạnh). Hỏi hình lục giác đó có bao nhiêu đỉnh?
\shortans{6}
\loigiai{
\textbf{Đáp số:} $6$\par
Hình lục giác có 6 cạnh và đúng $6$ đỉnh.
\textbf{Đáp số:} $6$.
}
\end{ex}

\begin{ex} %[1D2V08]
Xếp 4 que tính thành 1 hình thoi. Hỏi hình thoi đó có bao nhiêu cạnh?
\shortans{4}
\loigiai{
\textbf{Đáp số:} $4$\par
Hình thoi có đúng $4$ cạnh bằng nhau.
\textbf{Đáp số:} $4$.
}
\end{ex}
